\documentclass{article}
\usepackage{amsmath,amssymb}
\usepackage[vocab={geometry}]{leanify}
\newtheorem{lemma}{Lemma}
\begin{document}

\section{Notation}

We write $|PQ|$ for the distance between $P$ and $Q$.
We write $d(P,Q)$ for the distance between $P$ and $Q$.
We write $[ABC]$ for the area of the triangle with vertices $A$, $B$ and $C$.
We write $\alpha(A,B,C)$ for the angle at $A$ between $B$ and $C$.
We write $O$ for the origin.
We write $E$ for the point $(3,4)$.
We write $\Gamma(A,r)$ for the circle with centre $A$ and radius $r$.

\section{The plane}

\begin{lemma}\label{triangle}
Let $A$, $B$ and $C$ be points. Then $d(A,C)\le d(A,B)+d(B,C)$.
\end{lemma}

\begin{lemma}\label{midpoint}
Let $A$ and $B$ be points, and let $M$ be the midpoint of $A$ and $B$. Then $d(A,M)=d(M,B)$.
\end{lemma}

\begin{lemma}\label{coordinates}
$d(O,E)=5$, and $E$ is distinct from $O$.
\end{lemma}

\begin{lemma}\label{angles}
Let $A$, $B$ and $C$ be points with $[ABC]>0$. Then $\alpha(A,B,C)+\alpha(B,C,A)+\alpha(C,A,B)=\pi$.
\end{lemma}

\begin{lemma}\label{pythagoras}
Let $A$, $B$ and $C$ be points. If $\alpha(C,A,B)=\frac{\pi}{2}$, then $d(A,B)^2=d(A,C)^2+d(B,C)^2$.
\end{lemma}

\begin{lemma}\label{circle}
Let $A$ be a point. Every point $P$ lying on $\Gamma(A,2)$ is distinct from $A$.
\end{lemma}

\begin{lemma}\label{radius}
Let $A$ be a point. For every point $P$ lying on $\Gamma(A,2)$, $d(A,P)=2$.
\end{lemma}

\begin{lemma}\label{bisector}
Let $A$ and $B$ be points, and let $P$ be a point equidistant from $A$ and $B$. Then $[ABP]\ge0$.
\end{lemma}

\end{document}
